% 本文件由 LaTeX 排版 Agent 根据有效 translation.json 自动生成。
% author=Augustine of Hippo; work_id=1504; title=论人的义德成全

\chapter{论人的义德成全}

% structure: p0001 source=h1 target=omitted reason=page-title-already-rendered-by-parent

\Person{Augustine} 致其圣洁的弟兄及同为主教的 \Person{Eutropius} 与 \Person{Paulus}。\par

% structure: p0003 source=h2 target=section reason=relative-heading-rank; base=h2

\section{第1章}

二位的爱德，在你们两位身上是如此巨大且圣洁，以至于服从其命令乃是喜乐；这使我有义务回应一些据称是 \Person{Cœlestius} 所作的命题；因为你们给我的文稿的标题便是如此：\Quote{据称是 \Person{Cœlestius} 的命题}。至于这个标题，我认为它并非出自他本人，而是出自那些从 \Place{Sicily} 带来这作品的人——据说 \Person{Cœlestius} 并不在那里——但那里有许多人大言不惭地声称持有类似他的观点，并按宗徒的话说，\BibleQuote{他们自己受了欺骗，也误导了别人。} \BibleRef{弟后 3:13} 然而，这些观点确实是他或他的一些同路人的——我们也很可以相信。因为上述的简短命题，或者更确切地说，论断，与他的意见绝无不合之处——我在另一部确为他所著的作品中曾见到他如此表达。因此，我想那些带来这消息的弟兄们在 \Place{Sicily} 所听到的，即 \Person{Cœlestius} 教导或写下这类意见，是有充分理由的。我虽愿意，若可能的话，以如此兄弟般仁爱加诸于我的义务来回应，使我的答复也同样简短。但除非我同时也陈明我所答复的论点，否则谁能判断我答复的价值呢？然而，我将尽力而为，也赖天主的仁慈和你们的祈祷，使讨论不至冗长无益。\par

% structure: p0005 source=h2 target=section reason=relative-heading-rank; base=h2

\section{第2章}

% structure: p0006 source=h3 target=subsection reason=relative-heading-rank; base=h2

\subsection{\OriginalTerm{（1）Cœlestius 的第一条简论}{（1）The First Breviate of Cœlestius}}

他说：\Quote{首先，必须询问那否认人能不犯罪生活的人：每一种罪是什么？——是那种可以避免的罪？还是不可避免的？如果是不可避免的，那就不算罪；如果可以避免，那么人就可以避免那可以避免的罪而生活。没有理性或正义容许我们把无论如何不能避免的事称为罪。} 对此我们的答复是：罪是可以避免的，如果我们的败坏本性因天主的恩宠，借着我们的主耶稣基督得医治。因为，在它未健全的程度上，它要么因盲目而看不见，要么因软弱而做不到它所当行的；\BibleQuote{因为本性和心神彼此敌对，以致人不能行他所愿意的事。} \BibleRef{迦 5:17}\par

% structure: p0008 source=h3 target=subsection reason=relative-heading-rank; base=h2

\subsection{\OriginalTerm{（2）第二条简论}{（2）The Second Breviate}}

他说：\Quote{接着我们必须问：罪是出于意志，还是出于必然？如果是出于必然，就不是罪；如果是出于意志，就可以避免。} 我们如前答复；并且为了得医治，我们向那在圣咏中曾向祂呼求的祈祷：\BibleQuote{求祢从我的困境中领我出来。}\par

% structure: p0010 source=h3 target=subsection reason=relative-heading-rank; base=h2

\subsection{\OriginalTerm{（3）第三条简论}{（3）The Third Breviate}}

他说：\Quote{再次我们必须问：罪是自然的？还是偶然的？如果是自然的，就不是罪；如果是偶然的，就是可分离的；并且如果是可分离的，就可以避免；因为可以避免，人就能避免那可以避免的而生活。} 对此的答复是：罪不是自然的；但本性（尤其是在那败坏的状态中——我们由此在本性上是\BibleQuote{可怒之子} \BibleRef{弗 2:3} ）缺乏避免罪的意志决心，除非得蒙天主借着我们的主耶稣基督的恩宠扶助和医治。\par

% structure: p0012 source=h3 target=subsection reason=relative-heading-rank; base=h2

\subsection{\OriginalTerm{（4）第四条简论}{（4）The Fourth Breviate}}

他说：\Quote{再次我们必须问：罪是什么？——是一个行为，还是一个事物？如果它是一个事物，它必然有一个作者；如果它被认为有一个作者，那么除了天主之外，似乎要引入另一个事物之作者。但若说这是不敬的，我们就被迫承认：每一种罪是一个行为，而不是一个事物。因此，如果它是一个行为，正因它是一个行为，它就可以避免。} 我们的答复是：罪无疑被称为行为，并且如此，而不是一个事物。但同样在身体上，跛行出于同样理由是一个行为，而不是一个事物，因为脚本身，或身体，或由于脚受伤而跛行的人，是事物；但人仍不能避免跛行，除非他的脚被治好。同样的变化可能发生在内在的人身上，但那是借着天主的恩宠，通过我们的主耶稣基督。致使此人跛行的缺陷本身既不是脚，也不是身体，也不是人，甚至也不是跛行本身；因为当然在没有行走时就没有跛行，尽管每当试图行走时，那导致跛行的缺陷仍然存在。因此，让他问：这个缺陷该叫什么名字？——他愿意称之为事物、行为，还是事物中导致畸形行为出现的坏属性？同样，在内在的人身上，灵魂是事物，偷窃是行为，贪婪是缺陷，即灵魂变坏的属性，即使在它没有为满足其贪婪而行动时也是如此，甚至在它听到诫命\BibleQuote{不可贪恋} \BibleRef{出 20:17} 并谴责自己，却仍然贪心时。然而，借着信德，它接受更新；换句话说，它\BibleQuote{却一天比一天得医治} \BibleRef{格后 4:16} ——却唯独靠着天主通过我们的主耶稣基督的恩宠。\par

% structure: p0014 source=h2 target=section reason=relative-heading-rank; base=h2

\section{第3章}

% structure: p0015 source=h3 target=subsection reason=relative-heading-rank; base=h2

\subsection{\OriginalTerm{（5）第五条简论}{（5）The Fifth Breviate}}

\Quote{我们必须再次，}他说，\Quote{探问人是否应当无罪的。无疑他应当。若他应当，他就能；若他不能，则他不应当。人若不应当无罪，则他应当有罪——而若确定罪是应当的，它就不再是罪了。若说这样是荒谬的，我们就必须承认人应当无罪；并且显然他的义务不大于他的能力。}我们以之前答复中使用的相同比喻来回答。当我们看到一个跛子有机会治愈跛足时，我们当然有权说：\Quote{那个人不应当跛足；若他应当，他就能。}然而，他并非一愿意就立刻能行；只有在经过治疗、药物辅助其意愿之后方能如此。同样的事发生在内在的人身上，关乎罪——即其跛足——藉着那\Quote{来不是召义人，而是召罪人}的恩宠；\BibleRef{玛 9:13}因为\Quote{健康的人不需要医生，而只有生病的人才需要。}\BibleRef{玛 9:12}\par

% structure: p0017 source=h3 target=subsection reason=relative-heading-rank; base=h2

\subsection{(6) 第六摘要}

\Quote{此外，}他说，\Quote{我们必须探问人是否被命令要无罪；因为他要么不能，那样他就没有被命令；要么因为他被命令，所以他就能。因为为什么那完全不能做的事会被命令呢？}答案是，人被最明智地命令以正直的脚步行走，目的是当他发现自己甚至对此也无能为力时，他就能寻求为内在的人提供的医治罪的跛足的补赎——即天主的恩宠，藉着我们的主耶稣基督。\par

% structure: p0019 source=h3 target=subsection reason=relative-heading-rank; base=h2

\subsection{(7) 第七摘要}

\Quote{我们要提出的下一个问题，}他说，\Quote{是，天主是否愿意人无罪。无疑天主愿意；并且毫无疑问祂有能力。因为谁如此鲁莽，竟会怀疑那祂无疑所愿意的事是不可能的呢？}这是答案。若天主不愿意人无罪，祂就不会派遣祂的无罪的圣子来医治人的罪。这事发生在信者身上，他们日日更新，\BibleRef{格后 4:16}直到他们的义德成为完全，如同完全康复的健康。\par

% structure: p0021 source=h3 target=subsection reason=relative-heading-rank; base=h2

\subsection{(8) 第八摘要}

\Quote{再者，必须问这个问题，}他说，\Quote{天主愿意人怎样——有罪，还是无罪？无疑，祂不愿意他有罪。我们必须思量，若说人有权成为天主所不愿意的有罪的人，而否认他有权成为天主所愿意的无罪的人，这是何等不虔的亵渎；好像天主创造了任何人为了这样的结果——他能成为祂不要他成为的，而不能成为祂要他成为的；并且他应过一种违背祂旨意的生活，而不是与祂旨意相符的生活。}事实上这已经回答过了；但我看到我必须在此补充一点：我们是因望德而得救。\Quote{但所见到的望德不是望德；因为人看见了，为什么还盼望呢？但我们若盼望那未见到的，就要耐心等待。}\BibleRef{罗 8:24}因此，完全的义德只会在完全的健康达到时才达到；而这完全的健康将在爱德充满时来到，因为\Quote{爱是法律的满全；}\BibleRef{罗 13:10}而那时爱德的充满将来到，当我们\Quote{看见祂怎样，就怎样。}\BibleRef{若一 3:2}当信德达到所见之实象时，爱德再不能增添什么。\par

% structure: p0023 source=h2 target=section reason=relative-heading-rank; base=h2

\section{第四章}

% structure: p0024 source=h3 target=subsection reason=relative-heading-rank; base=h2

\subsection{(9) 第九摘要}

\Quote{我们需要解决的下一问题，}他说，\Quote{是：人怎样成为有罪的？——出于本性的必然，还是出于自由选择？若出于本性的必然，他无可指责；若出于自由选择，则问题来了：他从谁领受了这自由选择？无疑，从天主。然而，天主所赐的当然是好的。这是无可否认的。那么，若一样东西倾向于邪恶多于良善，它又怎能被证明是好的呢？因为若藉着它人能有罪而不能无罪，它就是倾向于邪恶多于良善了。}答案是：人之所以有罪是出于自由选择；但随后紧跟着惩罚性的败坏，从自由产生了必然。因此信德向天主呼喊：\Quote{求祢引领我脱离我的困境。}在如此困境中，我们要么无法理解我们所想要的，要么（虽有愿望）我们不够强壮去完成我们已经理解的。而正是这自由本身，由解放者应许给信者。祂说：\Quote{若圣子使你们自由，你们就真正自由了。}\BibleRef{若 8:38}因为本性被它因意志所陷入的罪征服，已经失去了自由。因此另一处圣经说：\Quote{因为人被谁制伏，就为谁所奴役。}\BibleRef{伯后 2:19}既然\Quote{健康的人不需要医生，而只有生病的人才需要；}\BibleRef{玛 9:12}同样，不是自由的人需要解放者，而是只有被奴役的人。因此向祂发出解放的欢呼：\Quote{祢从困境的急难中拯救了我的灵魂。}因为真正的自由也是真正的健康；这健康若意志保持良善，就永远不会失去。但因为意志犯了罪，有罪的艰难必然追随着罪人，直到他的软弱被完全医治，并重新获得这样的自由：一方面，持久地愿意幸福地生活；另一方面，自愿而幸福地必然过德性生活，永不犯罪。\par

% structure: p0026 source=h3 target=subsection reason=relative-heading-rank; base=h2

\subsection{(10) 第十摘要}

\Quote{既然天主造人是好的，}他说，\Quote{并且，除了造人好之外，还命令人做好事，那么，我们若认为人是恶的（因为他既不是被造成这样，也不是被命令这样），并否认他有能力成为善的（尽管他被造成这样，也被命令这样做），这是多么不敬啊！}我们的回答是：既然那时不是人自己，而是天主造人成为善的；同样，也是天主，而不是人自己，重新造人成为善的，同时把他从自己凭着意愿、相信并呼求这种解救而作的恶中解放出来。但这一切是由内在的人日渐更新\BibleRef{格后 4:16}，借着我们的主耶稣基督的天主恩宠而成就的，为的是外在的人在末日复活，进入不是惩罚的，而是生命的永生。\par

% structure: p0028 source=h2 target=section reason=relative-heading-rank; base=h2

\section{第五章}

% structure: p0029 source=h3 target=subsection reason=relative-heading-rank; base=h2

\subsection{（十一）第十一简述}

\Quote{必须提出的下一个问题是，}他说，\Quote{所有的罪是以多少种方式表现出来的？如果我没弄错，有两种：要么做了禁止的事，要么没有做命令的事。现在，正如一切禁止的事都能避免是肯定的，同样，一切命令的事都能做到也是肯定的。因为禁止或命令那些既不能避免也不能完成的事是徒劳的。既然我们不得不承认人既能避免一切禁止的事，又能做一切命令的事，我们怎能否认人有可能无罪呢？}我的回答是：在圣经中有许多神圣的诫命，全部提及太费力；但主，祂在世上完成了并缩短了祂的话语，明确宣告法律和先知系于两条诫命\BibleRef{玛 22:40}，好让我们明白，凡天主命令我们的其他一切，都归结于这两条诫命，并应归于它们：\Quote{你应全心、全灵、全意爱上主你的天主；}\BibleRef{玛 22:37}和\Quote{你应爱你的近人如你自己。}\BibleRef{玛 22:39}\Quote{在这两条诫命上，}祂说，\Quote{就挂全法律和先知。}\BibleRef{玛 22:40}因此，凡天主法律禁止我们的，以及凡命令我们做的，我们被禁止和命令，直接目的在于实现这两条诫命。或许一般的禁令是：\Quote{不可贪恋；}\BibleRef{出 20:27}一般的命令是：\Quote{你应爱。}\BibleRef{申 6:5}因此，宗徒保禄在某一处简短地包含了这两者，用这些话表达禁令：\Quote{不可与此世同化，}\BibleRef{罗 12:2}用这些话表达命令：\Quote{但要因心思的更新而变化。}\BibleRef{罗 12:2}前者属于消极诫命，即不贪恋；后者属于积极诫命，即爱。前者关乎节制，后者关乎正义。前者命令避开恶，后者命令追求善。藉着戒除贪恋，我们脱去旧人；藉着彰显爱，我们穿上新人。但除非天主赐予恩赐，没有人能节制\BibleRef{智 8:21}；天主的爱也不是由我们自己倾注在我们心中，而是由所赐给我们的圣神\BibleRef{罗 5:5}。然而，这在那些因愿意、相信和祈祷而进步的人中天天发生，他们\Quote{忘记背后，努力面前。}\BibleRef{斐 3:13}因为法律之所以灌输所有这些诫命，是为了当人未能履行它们时，他不要因骄傲而自高自大，而是在极其疲倦中投奔恩宠。这样，法律就完成了它作为\Quote{启蒙师}的职分，如此吓唬人，\Quote{引他归于基督}，使他爱祂\BibleRef{迦 3:24}。\par

% structure: p0031 source=h2 target=section reason=relative-heading-rank; base=h2

\section{第六章}

% structure: p0032 source=h3 target=subsection reason=relative-heading-rank; base=h2

\subsection{（十二）第十二简述}

\Quote{问题又来了，}他说，\Quote{人如何不能无罪——由于他的意志，还是由于本性？若是由于本性，那就不是罪；若是由于他的意志，那么意志很容易被意志改变。}我们回答，提醒他应反思，说——不仅仅是可能的（无疑，当天主恩宠辅助时这不可否认），而且意志被意志改变是\Quote{\Emph{很容易的}}，这是极端的傲慢；而宗徒却说：\Quote{血肉的欲望相反圣神，圣神的相反血肉；二者彼此敌对，以致你们不能做你们所愿意的事。}\BibleRef{迦 5:17}他没有说：\Quote{二者彼此敌对，以致你们不愿做你们所能的事，}而是\Quote{以致你们不能做你们所愿意的事。}那么，血肉的欲望——它当然是有罪的和败坏的，无非是犯罪的欲望，关于此，同一位宗徒教导我们不要让它\Quote{在你们必死的身体上为王}\BibleRef{罗 6:12}；藉此说法，他足够清楚地告诉我们，在我们的必死身体中必有一种存在，不得允许它在那里统治——这血肉的欲望，我说，怎么没有被意志改变呢？宗徒在\Quote{以致你们不能做你们所\Emph{愿意}的事}这些话语中清楚地暗示了意志存在，如果意志那么容易被意志改变的话？当然，我们并非藉此论证归咎于灵魂或身体的本性，这些是天主创造的，且全然是善的；而是说，它被自己的意志败坏后，没有天主的恩宠就不能被治好。\par

% structure: p0034 source=h3 target=subsection reason=relative-heading-rank; base=h2

\subsection{（十三）第十三简述}

\Quote{我们下一个要问的问题是，}他说，\Quote{如果人不能无罪，这是谁的过错——人自己的，还是别人的？若是人自己的，如果他不能成为他所不能成为的人，这如何是他的过错呢？}我们回答说，人之所以不能无罪，这是人的过错，因为正是由于人单独的意志，才导致他陷入了一种仅凭人单独的意志无法克服的必然性。\par

% structure: p0036 source=h3 target=subsection reason=relative-heading-rank; base=h2

\subsection{(14) 第十四摘要}

\Quote{再者，必须问这个问题：}他说，\Quote{如果人的本性是善的——除了马尔西翁或摩尼教徒，无人敢否认——那么若它不可能脱离邪恶，它又何以为善？因为谁能否认一切罪都是邪恶？}我们回答说：人的本性既是善的，也能脱离邪恶。因此我们热切祈祷：\BibleQuote{救我们脱离邪恶。}\BibleRef{玛 6:13}的确，这种解脱尚未完全实现，只要灵魂仍被那趋向腐朽的身体所压迫。\BibleRef{智 9:15}然而，这过程正借着恩宠，通过信德得以实现，以便不久就能说：\BibleQuote{死亡啊，你的争战在哪里？死亡啊，你的刺在哪里？死亡的刺就是罪，而罪的权势就是法律；}\BibleRef{格前 15:35}因为法律通过禁止罪恶，只会增加对它的欲望，除非圣神将那爱广传开来，那时我们将面对面地看见，那爱将是圆满而完美的。\par

% structure: p0038 source=h3 target=subsection reason=relative-heading-rank; base=h2

\subsection{(15) 第十五摘要}

\Quote{此外，还必须这样说：}他说：\Quote{天主确实是正义的；这不可否认。但天主将一切罪归于人。我想，这也必须承认：凡不被算为罪的，就不是罪。现在，如果有任何无可避免的罪，那么天主如何被称为正义的，当祂被认为将无可避免的罪归于任何人时？}我们回答说：很久以前就曾对骄傲之人宣告：\Quote{主不归罪于他的人是有福的。}祂不将罪归于那些怀着信德向祂说\Quote{\BibleQuote{宽免我们的罪债，如同我们也宽免我们的债务人}}的人。\BibleRef{玛 6:12}祂合理地省略了这种归算，因为祂所说的就是正义的：\Quote{\BibleQuote{你们用什么度量衡量，也要用什么度量衡量你们。}}\BibleRef{玛 7:2}然而，那罪就是：在其中要么没有应有的爱，要么爱少于应有的程度——无论它能否被人的意志所避免；因为当它能避免时，人现时的意志行了那罪；但如果它不能避免，则是人过去的意志行了它；然而，它还是可以被避免的——不过，不是在骄傲的意志受赞美时，而是在谦卑的意志得到帮助时。\par

% structure: p0040 source=h2 target=section reason=relative-heading-rank; base=h2

\section{第七章}

% structure: p0041 source=h3 target=subsection reason=relative-heading-rank; base=h2

\subsection{(16) 第十六摘要}

在所有这些辩论之后，这辩论的作者亲自以与另一人辩论的形象出现，并把自己表现为正受审查，且被他的审查者如此质问：\Quote{请指给我一个没有罪的人。}他回答说：\Quote{我指给你一个能没有罪的人。}他的审查者便对他说：\Quote{他是谁？}他回答说：\Quote{就是你自己。}\Quote{但是，}他补充说，\Quote{如果你说：\InnerQuote{我反正不能没有罪}，那么你必须回答我：\InnerQuote{那是谁的过错？}如果你接着说：\InnerQuote{是我自己的过错}，你又会被问：\InnerQuote{如果你不能没有罪，这怎么会是你的过错呢？}}他又把自己表现为受审查，并如此被质问：\Quote{你自己没有罪吗？你说一个人可以没有罪？}他回答说：\Quote{我没有罪，是谁的过错？但是，}他继续说，\Quote{如果他回答说：\InnerQuote{是你自己的过错}，那么答案将是：\InnerQuote{既然我不能没有罪，怎么会是我的过错呢？}}现在，我们对所有这些连续争论的回答是：他们之间本不应就这些言辞挑起任何争论；因为他从不敢肯定某个人（无论是别人还是他自己）没有罪，而只是回答说：他\Emph{能够}没有罪——这个立场我们自己并不否认。问题只在于：他何时能够，以及通过谁能够？如果是在当下，那么任何一个居于这死亡身体内的忠实灵魂都不应献上这样的祈祷，也不应说这样的话：\Quote{\BibleQuote{宽免我们的罪债，如同我们也宽免我们的债务人}}\BibleRef{玛 6:12}，因为在神圣的圣洗中，过去的一切罪债都已得到宽免。但是，若有人试图说服我们：这样的祈祷不适用于基督忠实的肢体，那他实际上无非是承认自己不是基督徒。再者，如果是人靠自己就能没有罪，那么基督就白白死了。但\Quote{基督并没有白白死去。}因此，人不可能没有罪，即使他愿意，除非他借着我们的主耶稣基督得到天主的恩宠之助。为使这圆满得以达成，在成长中的基督徒中甚至现在就已在进行一种训练，而在与死亡的斗争结束后，因信德和望德运作而珍爱的那爱，必在瞻仰和享有中圆满成全。\par

% structure: p0043 source=h2 target=section reason=relative-heading-rank; base=h2

\section{第八章}

% structure: p0044 source=h3 target=subsection reason=relative-heading-rank; base=h2

\subsection{(17) 离开身体是一回事，从这死亡的身体中被释放出来是另一回事}

接着，他试图以圣经的见证来确立他的论点。让我们仔细观察他所作的辩护。他说：\Quote{有些经文证明人受命成为无罪。}对此我们的回答是：问题根本不在于这样的命令是否被给予，因为事实足够清楚；而在于那明显被命令的事，在这个死亡的肉身中是否可能实现，在这肉身中，\BibleQuote{肉情相反圣神，圣神相反肉情，以致我们不能做我们所愿意的事。}\BibleRef{迦 5:17}并非每一个结束现世生命的人都从这死亡的肉身中得释放，只有那些在现世生命中领受了恩宠，并借着以善行度日证明自己没有白白领受恩宠的人，才能得释放。因为离开肉身是一回事（所有人在他们现世生命的最后一天都必须如此），而从这死亡的肉身中被释放是另一回事——这唯独是天主的恩宠，通过我们的主耶稣基督，赐予祂忠信的圣人们。确实，完美的赏报是在今生之后赐予的，但仅赐予那些在现世生命中已经获得了这种赏报之功劳的人。因为没有人能在离开之后达到义德的圆满，除非他在这里，借着饥渴慕义，跑完了他的赛程。\BibleQuote{饥渴慕义的人是有福的，因为他们要得饱饫。}\BibleRef{玛 5:6}\par

% structure: p0046 source=h3 target=subsection reason=relative-heading-rank; base=h2

\subsection{(18) 此世之义包含三部分——禁食、施舍与祈祷}

因此，只要我们\BibleQuote{离主而居，便凭信德往来，并非凭目睹；}\BibleRef{格后 5:6}因此经上说：\BibleQuote{义人因信德而生活。}\BibleRef{哈 2:4}我们在此世旅居中的义德就是：我们向着那完美的、圆满的义德迈进，在那义德中，在祂光荣的面前将会有完美而圆满的爱德；并且现在我们坚持我们道路的正直与完美，借着\BibleQuote{刻苦自己的身体，使它服从}\BibleRef{格前 9:27}，借着愉快而由衷地施行施舍，同时施与恩惠并宽恕得罪我们的人，并且借着\BibleQuote{恒心祈祷}\BibleRef{罗 12:12}——并以纯正的教义做这一切，其上建基于正确的信德、坚定的望德和纯洁的爱德。这就是我们现在的义德，我们在此走过我们的路程，饥渴慕义，盼望那完美的、圆满的义德，以便将来得以满足。因此，我们的主在福音中（说了\BibleQuote{你们应当小心，不要在人前施行你们的义德，为叫他们看见；}\BibleRef{玛 6:1}）为了叫我们不按人的光荣来衡量我们的生活道路，在阐述义德本身时宣告：除非有这三样——禁食、施舍、祈祷，就没有义德。在\Emph{禁食}中，祂表明了身体的完全征服；在\Emph{施舍}中，祂表明了意愿和行为的一切良善，无论是给予还是宽恕；在\Emph{祈祷}中，祂暗示了神圣渴望的一切规则。因此，虽然借着身体的征服，那种贪欲得到了抑制（这种贪欲不仅应被约束，而且应被彻底消除——在那种完美的义德状态中，它将完全不存在，因为罪将被绝对排除）——然而，即使在允许和正当的事物使用中，它也常常表现出过度的欲望。在那种真正的施惠中，义人顾及邻人的福祉，有时会做出有害的事情，尽管原以为是有益的。有时，由于软弱，当所花费的恩惠和劳苦要么不足以满足对象的需要，要么在环境中用处不大时，我们便会感到一种失望，这种失望玷污了那种\Quote{愉快}，这种愉快确保了\Quote{施与者}得到天主的赞许。\BibleRef{格后 9:7}然而，这种悲伤的痕迹，随着各人在善意上的进步程度而或大或小。如果这些考虑以及诸如此类的考虑被适当衡量，那么我们在祈祷中说\BibleQuote{宽免我们的罪债，如同我们也宽免得罪我们的人}\BibleRef{玛 6:12}就是正确的。但我们在祈祷中所说的，我们必须实行，甚至要爱我们的仇敌；或者，如果任何一个仍在基督内为婴孩的人尚未达到这一点，那么无论如何，每当一个得罪他的人悔改并请求宽恕时，他必须从心底里施行宽恕，如果他想让他的天父垂听他的祈祷。\par

% structure: p0048 source=h3 target=subsection reason=relative-heading-rank; base=h2

\subsection{(19) 爱的诫命将在来世得以完美成全}

在这祷词中，除非我们有意争辩，否则便有一面足够明亮的镜子放在我们眼前，从中可以看见义人的生活，他们因信德而生活，跑完他们的赛程，虽然他们并非没有罪过。因此他们说：\Quote{求你宽恕我们，}因为他们尚未到达赛程的终点。因此宗徒说：\Quote{这不是说我已经达到了，或已经是成全的了……弟兄们，我并不认为自己已经获得了；我只顾一件事：忘记背后的，努力前面的，向着目标奔驰，为得天主在基督耶稣内高天召叫的奖品。所以，我们凡是成全的人，都该怀有这种心情。}\BibleRef{斐 3:12} 换言之，我们凡是跑得端正的人，都该这样决心：既然尚未成全，就要沿着我们至今跑得端正的道路，向着成全奔跑，为的是\Quote{等到那圆满的来到，那局部的就要消逝；}\BibleRef{格前 13:10} 也就是，不再是局部的，而是成为整体和圆满。因为信德和望德将立刻被实质本身所取代，不再是被相信和盼望的，而是被看见和把握。然而爱德，三者中最大的，不是要被取代，而是被增长和满全——以全然的视觉默观它曾以信德所见到的，并在实际的享有中获得它曾只在望德中拥抱的。然后，在这爱德的圆满中，将满全这条诫命：\Quote{你当全心、全灵、全意爱上主你的天主。} 因为只要还存有任何肉欲的残余，需要用节制的缰绳加以约束，天主就绝不能被称为以全心爱着的。因为肉欲不是没有灵魂而贪欲；虽然说是肉欲在贪欲，因为灵魂在肉欲地贪欲。在那完美状态中，义人将毫无罪过地生活，因为在他肢体中不再有相反他心思之法律的律法，\BibleRef{罗 7:23} 但他将全心、全灵、全意爱天主，\BibleRef{玛 22:37} 这是第一条且最大的诫命。为什么这样的成全不该吩咐给人呢，尽管在今生可能没有人达到？因为如果我们不知道要跑向哪里，我们就不能正确地跑。但除非在诫命中指明，又怎能知道呢？因此，让我们\Quote{这样跑，好能得到奖品。}\BibleRef{格前 9:23} 因为所有跑得正确的都将得到——不像在剧场的竞赛中，所有人都跑，但只有一人赢得奖品。\BibleRef{格前 9:24} 让我们奔跑，怀着信德、望德、渴慕；让我们奔跑，制服身体，喜乐而热心地行哀矜——施予恩慈和宽恕伤害，祈祷我们的力量在奔跑中得到助佑；让我们如此听从催促我们达到完美的诫命，以致不忽略向着爱德的圆满奔跑。\par

% structure: p0050 source=h2 target=section reason=relative-heading-rank; base=h2

\section{第九章}

% structure: p0051 source=h3 target=subsection reason=relative-heading-rank; base=h2

\subsection{(20) 谁可被称为无瑕行走；大罪与小罪}

作了这些预备说明之后，让我们仔细审阅我们所答复的人所引用的段落，好像是我们自己引用了它们似的。\Quote{在申命纪中，\InnerQuote{你应在上主你的天主面前作一个齐全的人。}\BibleRef{申 18:13} 同一卷书中又说，\InnerQuote{在以色列子民中不应有一个不齐全的人。}\BibleRef{申 23:17} 同样，救主在福音中说：\InnerQuote{你们应当是成全的，如同你们的天父是成全的。}\BibleRef{玛 5:48} 宗徒在致格林多人后书中说：\InnerQuote{最后，弟兄们，你们要喜乐。要成全。}\BibleRef{格后 13:11} 他又致书哥罗森人说：\InnerQuote{我们劝诫各人，教导各人，以一切智慧，好把各人在基督内成全地呈献于天主。}\BibleRef{哥 1:28} 同样，致斐理伯人说：\InnerQuote{你们做一切事，总不可抱怨，也不可争论，好使你们成为无可指摘和纯洁的，成为天主无瑕的子女。}\BibleRef{斐 2:14} 致厄弗所人书中说：\InnerQuote{愿我们的主耶稣基督的天主和父受赞美！祂在天上，在基督内，以各种属神的祝福，祝福了我们，因为祂在创世以前，已在基督内拣选了我们，为使我们成为圣洁无瑕的。}\BibleRef{弗 1:3} 又在致哥罗森人另一处说：\InnerQuote{你们从前也与天主隔绝，并因恶行在心意上与祂为敌；可是现今天主在祂血肉的身体内，借着死亡，使你们与自己和好，好把你们呈献在祂面前，成为圣洁、无瑕和无可指摘的。}\BibleRef{哥 1:21} 同样，致厄弗所人说：\InnerQuote{祂为自己预备了一个光荣的教会，没有瑕疵，没有皱纹，或其他类似的缺陷；而成为圣洁无瑕的。}\BibleRef{弗 5:26} 致格林多人前书中说：\InnerQuote{你们当醒悟，行义，不要犯罪。}\BibleRef{格前 15:34} 伯多禄前书中又说：\InnerQuote{为此，你们要束起心灵的腰，谨慎自守，一心寄望于耶稣基督显现时所带给你们的恩宠；……你们应像顺从的子女，不要随从你们以往无知时的情欲，但应如那召叫你们的圣者一样，在一切行为上成为圣的，因为经上记载：\InnerQuote{\BibleRef{肋 19:2}你们应是圣的，因为我是圣的。}}\BibleRef{伯前 1:13} 同样，有福的达味也说：\InnerQuote{上主，谁能寄居你的帐幕？谁能安息在你的圣山？是那行为无瑕、实行正义的人。}又在另一处说：\InnerQuote{我在祂面前要成为无瑕的。}又说：\InnerQuote{品行无瑕、遵行上主法律的人，是有福的。}同样，撒罗满记载说：\InnerQuote{上主喜爱圣洁的心，所有无瑕的人都蒙祂悦纳。}}\BibleRef{箴 11:20} 这些章节中，有些是劝勉正在奔跑的人要跑得完全；另外一些则指向终点，使人奔跑时向着它前进。一个人被称为行为无瑕，并非因为他已抵达旅途终点，而是因为他以无可指摘的方式向着终点迈进，远离当受罚的大罪，同时不忽略以施舍来洁净那些小罪。因为我们所走的道路，即我们达到完美的途径，是靠洁净的祈祷来洁净的。所谓洁净的祈祷，就是我们真诚地说：\Quote{宽恕我们，如同我们宽恕别人一样。}\BibleRef{玛 6:12} 因此，只要不归咎于我们，我们就不受责备，我们就可以继续向完美迈进，无可指摘，简言之，无瑕；而当我们最终达到这种完美境界时，就会发现，绝对没有任何需要靠宽恕来洁净的东西。\par

% structure: p0053 source=h2 target=section reason=relative-heading-rank; base=h2

\section{第十章}

% structure: p0054 source=h3 target=subsection reason=relative-heading-rank; base=h2

\subsection{（21）天主的诫命为谁沉重，为谁不沉重。圣经为何说天主的诫命并不沉重。诫命是人自由意志的证据，祈祷是恩宠的证据}

他随后引用了一些章节，以证明天主的诫命并不沉重。然而，谁不知道，既然总诫命是爱（因为\BibleQuote{诫命的终极就是爱}\BibleRef{弟前 1:8}，且\BibleQuote{爱是法律的满全}\BibleRef{罗 13:10}），那么凡由爱而非恐惧所完成的事，便不沉重呢？然而，那些试图以恐惧遵守天主诫命的人，却被它们所压迫。\BibleQuote{但完美的爱驱除恐惧}\BibleRef{若一 4:18}；而且，就诫命的负担而言，爱不仅卸下其重压，甚至如同翅膀将它举起。然而，为使这爱得以拥有——甚至在这死亡的身体中尽可能地拥有——意志的决心毫无裨益，除非藉着我们的主耶稣基督，由天主的恩宠所助佑。因为必须一再重申，这爱\BibleQuote{倾注在我们心中}，不是由于我们自己，而是\BibleQuote{藉着所赐予我们的圣神}\BibleRef{罗 5:5}。圣经之所以坚持天主的诫命并不沉重这一真理，别无他故，而是为使觉得它们沉重的心灵明白，它尚未领受那使主的诫命成为所赞颂的那样——即温柔而愉快的——的资源；并使之以意志的呻吟祈祷，以获得这轻松恩赐。因为那说\Quote{愿我的心纯全无过}，以及\Quote{请照祢的言语引导我的脚步，不要让任何不义管辖我}，以及\BibleQuote{愿祢的旨意奉行在人间，如同在天上}\BibleRef{玛 6:10}，以及\BibleQuote{不要让我们陷于诱惑}\BibleRef{玛 6:13}，以及其他类似祈祷（列举太长）的人，实际上是在为遵守天主诫命的能力祈祷。的确，一方面，如果我们的意志在其中毫无作用，就不会给我们任何遵守诫命的命令；另一方面，如果我们的意志本身足够，就没有祈祷的余地。因此，天主的诫命被赞颂为不沉重，为的是让觉得它们沉重的人明白，他尚未领受那除去它们沉重性的恩赐；并且让他不认为自己真正在遵守它们，当他如此遵守以致于它们对他成为沉重时。因为\BibleQuote{天主喜爱乐捐的人}\BibleRef{格后 9:7}。然而，当一个人觉得天主的诫命沉重时，不要因绝望而崩溃；他当迫使自己去寻找、去祈求、去敲门。\par

% structure: p0056 source=h3 target=subsection reason=relative-heading-rank; base=h2

\subsection{（22）证明天主的诫命并不沉重的章节}

随后，他引用了那些将天主自己的诫命描述为不沉重的章节：现在我们留意它们的证词。\Quote{因为，}他说，\Quote{天主的诫命不仅不是不可能的，而且甚至不沉重。在《申命纪》中：\InnerQuote{上主你的天主必再转来，为你好而喜悦，像祂喜悦你的祖先一样，只要你听从上主你天主的声音，遵守写在本书上的诫命、典章和法律，一心一意归向上主你的天主。因为我今日所吩咐你的这诫命，并不沉重，离你也不远：不是在天空中，以致你问：谁能上到天上，为我们取得，使我们听见而遵行？也不是在海外，以致你问：谁能渡海，为我们取得，使我们听见而遵行？这话离你很近，就在你口里，在你心里，在你手中，为要遵行。}\BibleRef{申 30:9} 同样在福音中，主说：\InnerQuote{凡劳苦和负重担的，你们都到我这里来，我要使你们安息。你们背起我的轭，跟我学，因为我是良善心谦的：你们必要找得你们灵魂的安息。因为我的轭是柔和的，我的担子是轻松的。}\BibleRef{玛 11:28} 同样在圣若望的书信中写道：\InnerQuote{这就是爱天主，就是遵守祂的诫命；而祂的诫命并不沉重。}\BibleRef{若一 5:3}} 听到这些出自法律、福音和书信的证词，让我们被建立于那恩宠上，此恩宠是那些人不了解的，他们\BibleQuote{不明白天主的正义，企图建立自己的正义，而不服从天主的正义}\BibleRef{罗 10:3}。因为，如果他们不理解《申命纪》的那段话，如同宗徒保禄所引用的一样——即\BibleQuote{心里相信，可得以成义；口里承认，可得以得救}\BibleRef{罗 10:10}，因为\BibleQuote{健康的人不需要医生，而是有病的人需要}\BibleRef{玛 9:12}——他们当然（藉着他最后引用的若望宗徒的那段话：\BibleQuote{这就是爱天主，就是遵守祂的诫命；而祂的诫命并不沉重}\BibleRef{若一 5:3}）该受提醒：天主的诫命对于爱天主并不沉重，这爱唯独由圣神倾注在我们心中，而非人的意志抉择（若对此赋予过多，他们就不明白天主的正义）。然而，这爱将在所有对惩罚的恐惧被消除时，得以圆满。\par

% structure: p0058 source=h2 target=section reason=relative-heading-rank; base=h2

\section{第十一章}

% structure: p0059 source=h3 target=subsection reason=relative-heading-rank; base=h2

\subsection{（23）\Person{塞勒斯提乌斯}在面对公教徒所反对的圣经章节时，试图用其他章节来回避：第一段章节}

在这之后，他引用了那些通常被用来反对他们的章节。他并没有尝试解释这些章节，而是引用了看似相反的章节来论证，从而将问题缠绕得更紧。\Quote{因为，}他说，\Quote{有些圣经章节是反对那些凭圣经的权威而愚昧地以为自己能够摧毁自由意志或不犯罪的可能的人。因为，}他补充道，\Quote{他们惯于引用圣约伯的话来反对我们：\BibleQuote{谁能从污秽中取出洁净？没有一人；即使婴儿在世只活了一天。}}\BibleRef{约 14:4}然后他继续以其他引文对此作出某种回应；正如约伯自己所说：\Quote{虽然我是正义齐全的人，却成了被嘲笑的对象，}\BibleRef{约 12:4}——他不懂得人可以被称为义人，就是那在义德上已达到如此完善以至于非常接近它的人；而我们并不否认这在现世中当许多人凭信德行走时也能做到。\par

% structure: p0061 source=h3 target=subsection reason=relative-heading-rank; base=h2

\subsection{(24) 无罪与无过——有何不同}

同一件事也在另一段经文里被确认，他紧接着就引用了同一约伯的话：\BibleQuote{看，我离我的审判已经很近了，我知道我必被证明为义。}\BibleRef{约 13:18}这就是论到它时另一圣经章节所说的：\BibleQuote{祂必使你的义德如光照耀，你的审判如正午。}但他并没有说我已经到了那里，而是说\Quote{我离得很近}。的确，如果他所指的审判不是他自己将要行使的，而是他将在最后一日受审的审判，那么所有真诚祈祷的人：\BibleQuote{宽免我们的罪债，如同我们也宽免得罪我们的人，}\BibleRef{玛 6:12}都将在那样的审判中被证明为义。正是藉着这宽免，他们将被证明为义；因此，他们在此生所犯的一切罪，都因他们的爱德行为而被涂抹。因此主说：\BibleQuote{你们施舍吧；看，一切对你们都是洁净的。}\BibleRef{路 11:41}因为在末了，当义人将要进入所应许的国度时，将对他们说：\BibleQuote{我饿了，你们给了我吃，}\BibleRef{玛 25:35}等等。然而，无罪与无过是两回事：前者在今生只能归于独生子，而后者则可说及许多义人即使在今世；因为有一种良好生活的尺度，按照它，即使在此世的人际交往中也无法对他提出任何公义的控告。因为谁能公义地控告一个不愿任何人受害、并忠实地尽其所能善待所有人、从不怀有报复任何伤害他的人的愿望的人，使他能真诚地说：\Quote{如同我们宽免得罪我们的人}？但正因他真诚地说：\Quote{宽免我们，如同我们也宽免，}他显然承认自己并非无罪。\par

% structure: p0063 source=h3 target=subsection reason=relative-heading-rank; base=h2

\subsection{(25)}

因此这句话的力量在于：\BibleQuote{我的手没有不义，我的祈祷是纯洁的。}\BibleRef{约 16:18}因为他的祈祷纯洁是由于这个情况：当他真正宽恕别人时，在祈祷中祈求宽恕并无不当。\par

% structure: p0065 source=h3 target=subsection reason=relative-heading-rank; base=h2

\subsection{(26) 约伯为何如此受苦}

当他论及主说：\BibleQuote{祂无缘无故给了我许多伤痕，}\BibleRef{约 9:17}请注意他的话不是说祂给的伤痕\Emph{都有缘故}，而是\Quote{许多伤痕无缘无故}。因为不是因他众多的罪而给了他许多伤痕，而是为了考验他的耐心。的确，他承认自己在另一段经文中并非没有罪，但认为他应当受到的痛苦较少。\BibleRef{约 6:2}\par

% structure: p0067 source=h3 target=subsection reason=relative-heading-rank; base=h2

\subsection{(27) 谁可被称为遵守主的道路；何为偏离和离开主的道路}

接着，至于他所说：\BibleQuote{因为我遵守了祂的道路，没有偏离祂的诫命，我也不会离开它们，}\BibleRef{约 23:11}遵守天主的道路的人不会偏离以致放弃它们，而是奔跑在路上不断进步；虽然软弱时会偶尔绊倒或跌倒，但他仍然向前，越来越不犯罪，直到达到不再犯罪的全德状态。因为不遵守他的道路就不能以别的方式进步。的确，偏离这些道路并最终成为背道者的人，绝不是那种虽然有罪却永不止息地与之战斗直到到达死亡不再争战之家的人。好吧，在我们现在的争取中，我们穿上了义德，我们在此生凭信德生活——好像穿着胸甲一样穿上它。\BibleRef{弗 6:14}我们也承担审判；甚至当审判不利于我们时，我们把它转变成有利于我们；因为我们成为自己的控告者并定自己的罪：因此有经文说：\BibleQuote{义人一开始发言就控告自己。}\BibleRef{箴 18:17}因此他也说：\BibleQuote{我披上义德，又穿上审判如披外氅。}\BibleRef{约 29:14}我们的服饰现在无疑通常是用于战争的甲胄而非和平的衣裳，因为仍有私欲要制服；以后当最后一个敌人——死亡——被消灭时，情况将会不同，\BibleRef{格前 15:26}那时我们的义德将圆满完全，再也没有敌人来骚扰我们。\par

% structure: p0069 source=h3 target=subsection reason=relative-heading-rank; base=h2

\subsection{(28) 何时可说我们的心不责备我们；何时善得以成全}

此外，关于约伯的这些话语：\Quote{我的心在我一生中绝不会责备我，}\BibleRef{约 27:6} 我们指出，正是在我们现世的生活中（我们因信德度此生活），我们的心才不责备我们，倘若我们赖以相信以获得义德的同一信德没有疏忽责备我们的罪。基于这一原则，宗徒说：\Quote{我所愿意的善，我不去行；而我所不愿意的恶，我却去行。}\BibleRef{罗 7:15} 现在，避免私欲偏情是一件善事，而义人（他因信德而生活）愿意这善事；\BibleRef{哈 2:4} 但他仍然做他所憎恶的事，因为他有私欲偏情，尽管他\Quote{不随从自己的情欲；}\BibleRef{德 18:30} 如果他这样行了，那时他确实是自己行了这事，以至于屈服于、默许并顺从罪的欲望。于是他的心责备他，因为它责备他自己，而非责备那住在他内的罪。但每当他不让罪在他必死的身体上为王，去顺从它的情欲，\BibleRef{罗 6:12} 并且不将他的肢体献与不义作为罪的工具，\BibleRef{罗 6:13} 罪无疑仍存在于他的肢体中，但它并不为王，因为它的欲望未被顺从。因此，当他做他所不愿意的事——换言之，当他愿意没有情欲，却仍有情欲时——他同意法律是善的：\BibleRef{罗 7:16} 因为法律所愿意的，他也愿意；因为他的愿望是不放纵私欲偏情，而法律明确地说：\Quote{不可贪恋。}\BibleRef{出 20:17} 既然他愿意法律所愿意做的事，他无疑同意法律：但他仍有情欲，因为他不是没有罪的；然而，行那事的已不再是他自己，而是那住在他内的罪。因此，\Quote{他的心在他一生中不责备他；}也就是说，在他的信德中，因为义人因信德而生活，以致他的信德就是他的生命。他确实知道，在他自己内——甚至在他的肉身内，那是罪的居所——没有善住着。然而，由于不向它屈服，他因信德而生活，并以此呼求天主帮助他对抗罪的斗争。此外，他愿意自己身上丝毫没有罪，但如何成全这善却不在于他。他所欠缺的不仅仅是\Quote{行}一件善事，更是它的\Quote{成全}。因为在他不屈服这一点上，他行善；在他憎恨自己的情欲这一点上，他再次行善；在他不停施舍这一点上，他也行善；在他宽恕得罪自己的人这一点上，他行善；在他祈求自己的过犯得宽恕——在他的祈祷中真诚宣告他也宽恕得罪自己的人，并祈求不受诱惑，脱离凶恶——这一点上，他行善。但如何成全这善却不在于他；然而，在最终的状态，当住在他肢体中的私欲偏情不复存在时，这成全就会实现。因此，他的心在责备那住在他肢体中的罪时，并不责备他；它也不能责备他不信。所以，\Quote{在他一生中}——也就是说，在他的信德中——他既未受自己内心的责备，也未被证明并非没有罪。约伯自己也承认这一点，他说：\Quote{我的罪没有一件能逃过你；你把我的过犯封在袋里，并标记了我是否无意中行了不义。}\BibleRef{约 14:16} 那么，关于他从\BibleBook{约伯}{传}中所引的段落，我们已尽力说明了它们应按何种意义来理解。然而，他未能解释他自己从同一\BibleBook{约伯}{传}中所引的话的含义：\Quote{谁能从污秽中纯洁起来？没有一人，即使他仅在地上为一天的婴儿。}\BibleRef{约 14:4}\par

% structure: p0071 source=h2 target=section reason=relative-heading-rank; base=h2

\section{第十二章}

% structure: p0072 source=h3 target=subsection reason=relative-heading-rank; base=h2

\subsection{(29) 第二段经文：谁可被称为远离一切恶事}

\Quote{他们惯于接下来引用，}他说，\Quote{这段话：\InnerQuote{人人都是说谎者。}}但这里他再次没有解决那些被他本人用来反驳自己的话；他所做的只是在那些不熟悉圣经的人面前提及其他明显相反的段落，从而使天主的话陷入冲突。他是这样说的：\Quote{我们回答他们说：在\BibleBook{}{户籍纪}中写道：\InnerQuote{人是真诚的}；而对于圣约伯，我们读到这样的赞辞：\InnerQuote{在乌兹地有一个人，名叫约伯；那人真诚、无瑕、正义、敬畏天主，远离一切恶事。}}\BibleRef{约 1:1} 我感到惊讶，他引用了这段话说约伯\Quote{远离一切恶事}，并希望它意指\Quote{远离一切罪；}因为他已经论证过罪不是一件事，而是一个行为。让他回想一下，即使它是一个行为，它仍然可以被称为一件事。然而，那远离一切恶事的人，要么从不向那常与他同在的罪屈服，要么即使有时被它紧紧抓住，也绝不被它压制；正如摔跤冠军，尽管有时被猛烈缠住，却不会因此失去使他成为更强者的能力。的确，我们读到过无可指责的人，不受控告的人；但我们从未读到过无罪的人，除了那一位人子，他也是天主的独生子。\par

% structure: p0074 source=h3 target=subsection reason=relative-heading-rank; base=h2

\subsection{(30) \Quote{人人都是说谎者}只由于他自己；但\Quote{人人都是真诚者}唯赖天主的恩宠}

\Quote{此外，在约伯本人中记载：\BibleQuote{他保持了真人的奇迹。}\BibleRef{约 17:8} 我们又在撒罗满的智慧书中读到：\BibleQuote{说谎的人不能记念她，但真理的人必在她内被找到。}\BibleRef{德 15:8} 又在默示录中：\BibleQuote{在他们口中找不出诡诈，因为他们是没有瑕疵的。}}\BibleRef{默 14:5} 对于所有这些陈述，我们回答并提醒我们的对手：一个人如何能被称作真实的，乃是藉着天主的恩宠和真理，而他自己无疑是个说谎者。因此经上说：\BibleQuote{每个人都撒谎。} 至于他所引用的关于智慧的那段经文，当经上说\BibleQuote{真理的人将在她内被找到}时，我们必须注意，人们被找到是说谎者，无疑地不是\Quote{\Emph{在她内}}，而是\Emph{在自身内}。正如在另一段经文：\BibleQuote{你们从前是黑暗，但如今在主内是光明}\BibleRef{弗 5:8}——当他说\BibleQuote{你们是黑暗}时，他没有加上\BibleQuote{在主内}；但在说了\BibleQuote{你们如今是光明}之后，他特意加上了\BibleQuote{在主内}，因为他们不可能在自身内是\BibleQuote{光明}；为叫\BibleQuote{那夸耀的，要在主内夸耀}。\BibleRef{格前 1:31} 实际上，默示录中的\Quote{无瑕疵者}之所以如此称呼，是因为\BibleQuote{在他们口中找不出诡诈。}\BibleRef{默 14:5} 他们并没有说他们没有罪：如果他们这样说，他们就是自欺，真理就不在他们内；\BibleRef{若一 1:8} 如果真理不在他们内，他们口中就会找出诡诈和不真实。然而，如果他们为了避免嫉妒，说他们不是没有罪，尽管他们是无罪的，那么这种不真诚本身就是谎言，对他们的描述也就不真实：\BibleQuote{在他们口中找不出诡诈。} 因此他们确实是\Quote{无瑕疵的}；因为他们宽恕了那些得罪他们的人，同样他们也因天主对自己的宽恕而被净化。现在请注意，我们已尽最大努力解释了他为自己所引用的引文应当如何理解。但他本人却完全没有解释\BibleQuote{每个人都撒谎}这段经文该如何解释；若不纠正那个使他相信人不需要天主恩宠的帮助、仅凭自己的自由意志就能成为真实的错误，他不可能做出解释。\par

% structure: p0076 source=h2 target=section reason=relative-heading-rank; base=h2

\section{第十三章}

% structure: p0077 source=h3 target=subsection reason=relative-heading-rank; base=h2

\subsection{（31）第三段经文。离开一切罪与已离开一切罪是两回事。\Quote{没有行善的}——这是指着谁说的}

他同样提出了另一个问题，正如我们将要展示的，但没有解决它；不仅如此，他甚至使问题更加困难，首先陈述了曾用来反对他的证据：\Quote{没有行善的，连一个也没有；} 然后诉诸看似相反的经文，以表明有行善的人。他无疑做到了这一点。然而，一个人不做善事与一个人不是没有罪是两回事，尽管他同时可能做许多善事。因此，他所引用的经文并不真正与\Quote{没有人在今生是没有罪的}这一陈述相矛盾。他自己并没有解释，在何种意义上经上说：\Quote{没有行善的，连一个也没有。} 这是他的话：\Quote{圣达味确实说：\BibleQuote{你要信赖主，并要行善。}} 但这是一条诫命，而不是已经完成的事实；这样的诫命，那些被论及\Quote{没有行善的，连一个也没有}的人，从未遵守。他补充说：\Quote{圣多俾亚也说：\BibleQuote{我儿，不要害怕我们遭受贫穷；如果我们敬畏天主，远离一切罪，并行善，我们必有许多福分。}} \BibleRef{多 4:21} 确实，人若已远离一切罪，必有许多福分。那时，没有灾祸会临到他；他也不需要祈祷说：\BibleQuote{救我们免于凶恶。} \BibleRef{玛 6:13} 尽管即使是现在，每个进步的人，以正直的目的不断前进，都在远离一切罪，当他愈接近正义状态的圆满和成全时，就愈远离罪；因为甚至私欲本身，即居住在我们肉体中的罪，在那些正在进步的人中也不会停止减少，尽管它仍然留在他们有死的肢体中。因此，远离一切罪——这是一个正在进行的进程——与已远离一切罪——这将在未来的成全状态中实现——是两回事。然而，即使是已经远离了恶并继续这样做的人，也必须被承认为行善者。那么，在他所引用且未解决的经文中，如何说\Quote{没有行善的，连一个也没有}呢？除非圣咏作者在那里斥责某个民族，在他们中没有一个人行善，他们愿意继续作\Quote{人子}，而不作天主的儿子；人藉着天主的恩宠成为善，以便行善。因为我们必须认为圣咏作者在这里指的是他在上下文中描述的\Quote{善}，他说：\BibleQuote{天主从天俯视人子，要看有没有明悟并寻求天主的。} 这样的善，即寻求天主，没有一个人追求它，连一个也没有；但这只限于那预定受灭亡的人群。天主正是用祂的先见俯视了这样的人，并宣判了。\par

% structure: p0079 source=h2 target=section reason=relative-heading-rank; base=h2

\section{第十四章}

% structure: p0080 source=h3 target=subsection reason=relative-heading-rank; base=h2

\subsection{（32）第四段经文。惟独天主是善的，意义何在。对天主而言，是善与是自身是同一回事}

\Quote{他们也同样,}他说，\Quote{引用救主的话：\InnerQuote{\BibleQuote{你为什么称我是良善的？除了天主一位外，没有良善的。}}} \BibleRef{路 18:19} 然而，他对此言论根本不加解释；他所做的只是提出其他一些似乎与之矛盾的不同段落，并引用它们来证明人也是良善的。以下是他所说的话：\Quote{我们必须用另一段经文来回答这段经文，在这段经文中，同一主说：\InnerQuote{\BibleQuote{善人从他心里所存的善，就发出善来。}} \BibleRef{玛 12:35} 又说：\InnerQuote{\BibleQuote{祂使太阳升起，光照善人，也光照恶人。}} \BibleRef{玛 5:45} 另一处又记载：\InnerQuote{\BibleQuote{善的事物是从起初创造的，}} \BibleRef{德 39:25} 再一处：\InnerQuote{\BibleQuote{善人必居住在地上。}} \BibleRef{箴 2:21}} 现在我们必须回答这一切：所讨论的这些段落必须与前面那段经文作同样的理解：\BibleQuote{除了天主一位外，没有良善的。} 或者因为一切受造之物，虽然天主造它们时非常好，但与造物主相比，便不是善的，因为实际上无法与祂相比。因为祂以一种超越而又非常恰当的方式论及自己说：\BibleQuote{\Emph{我是自有者。}} \BibleRef{出 3:14} 因此，我们眼前的这句话：\BibleQuote{除了天主一位外，没有良善的，} 与论及若翰的话有类似用法：\BibleQuote{他不是那光，} \BibleRef{若 1:8} 虽然主称他为\BibleQuote{灯，} 正如祂对门徒所说：\BibleQuote{你们是世界的光……人点灯，不放在斗底下。} \BibleRef{玛 5:14} 然而，与那光是\BibleQuote{真光，照亮每个来到世界上的人} \BibleRef{若 1:9} 相比，他不是光。或者，因为即使是天主的子女，当他们与自己将来在永恒完满中所成为的相比，他们既是善的，却仍然也是恶的。虽然我不敢对他们这样说（因为谁敢称那些以天主为父的人为恶呢？）除非主亲自说过：\BibleQuote{你们虽然邪恶，尚且知道把好东西给你们的儿女，何况你们的天父，岂不更将好的赐予求祂的人？} \BibleRef{玛 7:11} 当然，祂用\BibleQuote{你们的父}这一称呼，证明了他们已经是天主的子女；但同时祂毫不迟疑地说他们是\BibleQuote{邪恶的。} 然而，你那位作者并没有向我们解释，他们如何是善的，同时又\BibleQuote{除了天主一位外，没有良善的。} 因此，那个问\BibleQuote{我该做什么善事} \BibleRef{玛 19:16}的人，被劝勉去寻求祂 \BibleRef{路 10:27}，藉祂的恩宠才能成为善的；对祂而言，\Emph{是善的}无非就是\Emph{祂自己}，因为祂是不变之善，绝不可能邪恶。\par

% structure: p0082 source=h3 target=subsection reason=relative-heading-rank; base=h2

\subsection{(33) 第五段经文。}

\Quote{这段经文，}他说，\Quote{是他们的另一段引文：\InnerQuote{\BibleQuote{谁能说，我有一颗纯洁的心？}}} \BibleRef{箴 20:9} 然后他用几段经文回答，想要证明人可以有纯洁的心。但他没有告诉我们，他报告的那段被他用来反驳自己的经文应该如何理解，以免圣经在这段经文和他用来回答的经文中显得自相矛盾。我们这一方面确实回答他：\BibleQuote{谁能说，我有一颗纯洁的心？} 这句是前文\BibleQuote{当义王坐在宝座上时} \BibleRef{箴 20:8} 的恰当续句。因为无论一个人的义德多么伟大，当那位义王坐在祂的宝座上时，他应当反省和思考，以免发现某种应受责备的事，这些事确实逃过了他自己的注意。在那位义王的洞察之下，没有罪能逃脱，就连那些关于\Quote{谁能了解自己的过犯？} 所说的罪也不能。\BibleQuote{因此，当义王坐在祂的宝座上时……谁能说，他有一颗纯洁的心？或谁敢大胆地说，他纯洁无罪？} \BibleRef{箴 20:8} 除非是那些想要夸耀自己义德的人，而不是以审判者本身的仁慈为荣。\par

% structure: p0084 source=h2 target=section reason=relative-heading-rank; base=h2

\section{第十五章}

% structure: p0085 source=h3 target=subsection reason=relative-heading-rank; base=h2

\subsection{(34) 对立的经文}

然而，他接着引用来作为回答的经文是真实的，说：\Quote{救主在福音中宣称：\InnerQuote{心里纯洁的人是有福的，因为他们要看见天主。}\BibleRef{玛 5:8}达味也说：\InnerQuote{谁能登上主的山？谁能站在祂的圣所？那手洁心清的人。}又在另一处说：\InnerQuote{上主，求祢善待那些善良和心正的人。}同样在撒罗满的书中：\InnerQuote{财富对于良心无罪的人是有益的；}\BibleRef{德 13:24}接着在同一书中：\InnerQuote{远离罪恶，洁净你的手，从心里清除一切邪恶。}\BibleRef{德 38:10}同样在若望的书信中：\InnerQuote{如果我们心中不自责，我们就在天主前有坦然无惧的心，并且我们无论求什么，必从祂那里获得。}\BibleRef{若一 3:21}}因为这一切都是由意志、藉着信德、望德和爱德的操练，藉着制服身体，藉着施舍，藉着宽恕冒犯，藉着恳切祈祷，藉着祈求力量以在我们路程上前进，藉着真诚地说：\Quote{\InnerQuote{求祢宽恕我们，如同我们也宽恕他人，}以及\InnerQuote{不要让我们陷于诱惑，但救我们免于凶恶。}\BibleRef{玛 6:12}}藉着这个过程，确实实现我们的心得到洁净，所有的罪都被除去；当那公义的君王坐在祂的宝座上时，祂将因祂的仁慈赦免心中尚存而未洁净的一切罪，使整个人变得健全和洁净，得以看见天主。因为\Quote{\InnerQuote{不怜悯人的，要受无怜悯的审判；然而怜悯胜过审判。}\BibleRef{雅 2:13}}若非如此，我们之中还有谁有指望呢？\Quote{当公义的君王坐在祂的宝座上时，谁能夸口自己有一颗纯洁的心，或者谁能大胆地说自己无罪呢？}然而，藉着祂的仁慈，义人届时将完全而彻底地被洁净，像光辉的太阳在他们父的国度里照耀。\BibleRef{玛 13:43}\newline{}\par

% structure: p0087 source=h3 target=subsection reason=relative-heading-rank; base=h2

\subsection{（35）复活后教会将毫无瑕疵和皱纹}

那时，教会将完全而圆满地实现\Quote{\Quote{没有瑕疵、皱纹或任何这类东西}\BibleRef{弗 5:27}}的状况，因为那时它也要真实地成为光荣的。既然他加上了\Quote{光荣}这一形容词，说：\Quote{为使祂把教会呈现在自己面前，没有瑕疵、皱纹或任何这类东西}，他便充分表明了教会何时将没有瑕疵、皱纹或任何这一类东西——当然是在它成为光荣的时候。因为与其说当教会处于如此多的邪恶之中，或处于如此多的罪过之中，以及如此多大坏人的混杂之中，并遭遇不敬神之人的沉重责难时，我们应当说它是光荣的——因为君王事奉它（这一事实只产生更危险更痛苦的诱惑）——不如说它那时将是光荣的，当那使徒所说的事件发生时：\Quote{\Quote{当基督——你们的生命——显现时，那时你们也要与祂一同在光荣中显现。}\BibleRef{哥 3:4}}因为主自己，在祂取了奴仆的形体，藉以作为中保与教会结合时，若非藉着复活的荣耀，也没有被光荣化（因此说：\Quote{\Quote{圣神还未赐下，因为耶稣尚未受到光荣}\BibleRef{若 7:39}}），那么祂的教会怎能在复活之前就被描述为\Emph{光荣的}呢？因此，现在祂藉着\Quote{\Quote{水洗，藉着圣言}\BibleRef{弗 5:26}}洁净它，洗去它过去的罪，驱逐恶天使的权势；但那时，藉着使它所有健康的机能达于圆满，祂使它适于那光荣的状态，在那里它要毫无瑕疵或皱纹地照耀。因为\Quote{\Quote{祂所预定的，祂也召叫了；祂所召叫的，祂也成义了；祂所成义的，祂也光荣了。}\BibleRef{罗 8:30}}我猜想，正是在这奥迹下，说了这话：\Quote{\Quote{看，我今天和明天驱逐魔鬼，施行医治，第三天我便完成}或\Quote{得以成全。}\BibleRef{路 13:32}}因为祂是以祂身体的位格——这身体就是祂的教会——说这话的，以\Emph{\Quote{日子}}表示明确指定的时期，这也在祂复活时所预示的\Quote{第三天}中指明。\newline{}\par

% structure: p0089 source=h3 target=subsection reason=relative-heading-rank; base=h2

\subsection{（36）心正者与心洁者的区别}

我想，心灵正直的人与心灵洁净的人也有所不同。一个人心灵正直，是指他\Quote{忘记背后，努力面前的}\BibleRef{斐 3:13}，以正直的道路——即以正直的信德与意向——抵达那能使他心灵洁净而纯洁居住的完美境界。因此，在圣咏中，这些条件应当分别赋予不同的品格，正如经上所说：\Quote{谁能登上主的山？谁能站在祂的圣所？就是手洁心清的人。}那手洁的人将登上，那心清的人将站立——前者是现世的操作，后者是圆满的实现。对于这些，更应当理解为所记载的：\Quote{良心无罪的人，财富才是好的。}\BibleRef{德 13:24}那时，当一切贫穷——即一切软弱——都被除去之后，真正的财富或真福才会来临。人现在确实可以\Quote{远离罪恶}，因为他在前进的进程中离开罪恶，日渐更新；他可以\Quote{安排他的手}，指引他从事慈悲的工作，并\Quote{洁净他的心脱离一切邪恶}\BibleRef{德 38:10}——他也可以如此慈悲，以致所剩余的可凭白白的宽恕获得赦免。这确实是圣若望所说的话的健全而恰当的含义，毫无虚妄和空洞的夸耀：\Quote{如果我们的心不责备我们，我们在天主前就可以放心安泰；我们无论求什么，必由祂获得。}\BibleRef{若一 3:21}他在这一段话中明显向我们提出的警告，是要提防我们的心在祈祷和恳求时责备我们；也就是说，当我们恰好用这祈祷文说\Quote{求祢宽恕我们，如同我们宽恕别人一样}时，我们却因没有实行所说的话而良心不安，甚至失去勇气说出我们所未能做到的，从而丧失信实和恳切祈祷的信心。\par

% structure: p0091 source=h2 target=section reason=relative-heading-rank; base=h2

\section{第十六章}

% structure: p0092 source=h3 target=subsection reason=relative-heading-rank; base=h2

\subsection{(37) 第六处经文}

他还引用了这段圣经，这段经文常被用来反对他的团体：\Quote{世上没有一个行善而不犯罪的义人。}\BibleRef{训 7:20}他假装用其他经文来回答它——比如：\Quote{主论及圣约伯说：\InnerQuote{你曾用心察看我的仆人约伯没有？地上没有人像他那样十全十美，正直、敬畏天主、远离邪恶。}}\BibleRef{约 1:8}关于这段经文，我们已经作过一些评论。但他甚至没有试图向我们说明，一方面，如果这些话具有那样的含义，约伯在地上如何绝对无罪；另一方面，他所承认的圣经的话\Quote{世上没有一个行善而不犯罪的义人}\BibleRef{训 7:20}又如何能成立。\par

% structure: p0094 source=h2 target=section reason=relative-heading-rank; base=h2

\section{第十七章}

% structure: p0095 source=h3 target=subsection reason=relative-heading-rank; base=h2

\subsection{(38) 第七处经文。谁可称为无瑕疵的。为何在天主眼中没有人能成义}

他说：\Quote{他们还引用经文：\InnerQuote{因为在祢面前，凡有血肉的，没有一个能成义。}}他对此经文的矫饰回答，无非是显示圣经经文之间似乎相互矛盾，而我们的责任却是证明它们的一致性。他的原话是：\Quote{我们必须用福音论及圣匝加利亚和依撒伯尔的见证来反驳他们，因为福音说：\InnerQuote{二人在天主前都是义人，遵行主的一切诫命和法令，无可指摘。}}\BibleRef{路 1:6}当然，这两位义人曾在这诫命中读到洁净自己罪恶的方法。因为按照《希伯来书》论及\Quote{每一位大司祭是从人间选出来的}\BibleRef{希 5:1}所说，匝加利亚无疑也曾为自己的罪献祭。然而，应用于他的\Quote{无可指摘}一词的含义，我想我们已经充分解释过了。他接着说：\Quote{蒙福的宗徒说：\InnerQuote{使我们成为圣洁无瑕疵的。}}\BibleRef{弗 1:4}照他看来，这话是说我们应当这样，如果\Quote{无可指摘}是指那些完全没有罪的人。然而，如果\Quote{无可指摘}是指那些没有责备或非难的人，那么我们不可能否认即使在今世也有这样的人，并且仍然有；因为一个人没有罪并不等于他没有被控告的污点。因此，宗徒在选择按立的执事时，并没有说\Quote{若有人没有罪}，因为他找不到这样的人；而是说\Quote{若有人无可指责}\BibleRef{铎 1:6}，因为这样的人他当然能找到。但我们的对手并没有告诉我们，按照他的观点，我们应当如何理解圣经：\Quote{因为在祢面前，凡有血肉的，没有一个能成义。}这句话的含义很清楚，前面的句子更使其明朗：\Quote{求祢不要审问祢的仆人，因为凡有血肉的，在祢面前没有一个能成义。}他害怕的是审判，因此他渴望那胜过审判的仁慈。\BibleRef{雅 2:13}因为\Quote{求祢不要审问祢的仆人}这一祈祷的意思是：\Quote{不要按照祢自己来审判我，因为祢是无罪的；因为凡有血肉的，在祢面前没有一个能成义。}这毫无疑问是指现世生活而言，而\Quote{不能成义}所指向的是不属于这生命的完美义德境界。\par

% structure: p0097 source=h2 target=section reason=relative-heading-rank; base=h2

\section{第十八章}

% structure: p0098 source=h3 target=subsection reason=relative-heading-rank; base=h2

\subsection{(39) 第八处经文。论及由天主所生者不犯罪是何意义。犯罪者不能看见或认识天主是何意义}

\Quote{他们也引用}，他说，这段经文，\BibleQuote{如果我们说我们没有罪，就是欺骗自己，真理就不在我们内。}\BibleRef{若一 1:8} 他试图用明显矛盾的经文来反驳这个非常清楚的见证，这样说：\Quote{同一圣若望在这封书信中说：\InnerQuote{不过，弟兄们，我说这话，是要你们不犯罪。凡由天主生的，就不犯罪；因为天主的种子存留在他内，他不能犯罪。}\BibleRef{若一 3:9}\InnerQuote{凡由天主生的，就不犯罪；因为天主生他的保守他，那恶者不能触摸他。}\BibleRef{若一 5:18} 又在另一处，当论及救主时，他说：\InnerQuote{祂曾显现，为除去罪恶，凡住在祂内的，就不犯罪；凡犯罪的，没有看见祂，也没有认识祂。}\BibleRef{若一 3:5} 又一次：\InnerQuote{亲爱的，现在我们是天主的子女，将来如何，还未显明；但我们知道，当祂显现时，我们将相似祂，因为我们要看见祂实在怎样。凡对祂怀着这希望的，就洁净自己，如同祂是洁净的。}\BibleRef{若一 3:2}} 然而，尽管这一切经文都是真实的，他所引用的那一段也是真实的，但他没有提供任何解释：\BibleQuote{如果我们说我们没有罪，就是欺骗自己，真理就不在我们内。}\BibleRef{若一 1:8} 由此可知，就我们由天主而生，住在祂内——祂显现是为除去罪恶，即住在基督内——而言，我们不犯罪——这无非是\Quote{我们内心的人日日更新}；\BibleRef{格后 4:16} 但就我们由那人而生——\Quote{通过他，罪恶进入世界，死亡因罪而来，于是死亡传给了众人}\BibleRef{罗 5:12} ——而言，我们不是没有罪的，因为我们尚未脱离他的软弱，直到那日日更新的更新（因为这是我们由天主而生的根据）完全修复我们在第一个人类那里所出生的软弱，在这软弱中我们不是没有罪的。只要这软弱在我们内心的人中还存留，无论它在前进的人中如何日渐减少，\Quote{如果我们说我们没有罪，就是欺骗自己，真理就不在我们内}。然而，\BibleQuote{凡犯罪的，没有看见祂，也没有认识祂}\BibleRef{若一 3:6} 这句话虽然真实，因为那将在实际见中实现的看见和认识，在此生无人能看见或认识祂；但凭信德的看见和认识，有许多人犯罪——甚至背教者——他们曾经相信过祂；因此，根据尚属信德阶段的看见和认识，不能说他们没有看见祂或认识祂。但我认为应该这样理解：那有待成全的更新看见了祂并认识了祂；而那注定要毁坏的软弱既看不见祂也不认识祂。正是由于这软弱的残余（无论多少）仍存留于我们内心的人中，\Quote{当我们说我们没有罪时，我们欺骗自己，真理就不在我们内}。所以，虽然因着更新的恩宠\Quote{我们是天主的子女}，但由于我们内软弱的残余，\Quote{将来如何，还未显明；我们只知道，当祂显现时，我们将相似祂，因为我们要看见祂实在怎样}。那时不再有罪，因为内外的软弱都不再存留。\Quote{凡对祂怀着这希望的，就洁净自己，如同祂是洁净的}——洁净自己，当然不是靠自己，而是信靠祂并呼求那圣化祂圣人的主；这圣化最终完成时（目前只是在日渐进步和增长），将永远除去我们一切软弱的残余。\par

% structure: p0100 source=h2 target=section reason=relative-heading-rank; base=h2

\section{第十九章}

% structure: p0101 source=h3 target=subsection reason=relative-heading-rank; base=h2

\subsection{（40）第九段经文}

\Quote{这段经文，}他说，\Quote{是他们引用的：\BibleQuote{不靠那愿意的，也不靠那奔跑的，只靠那施以仁慈的天主。}}\BibleRef{罗 9:16}他观察到，回答他们的话来自同一位宗徒在另一处的话：\BibleQuote{让他做他所愿意的。}\BibleRef{格前 7:36}他又引用了\BibleBook{}{费肋孟书}的另一段话，论及敖泥西摩，［圣保禄］说：\Quote{我本想把他留在身边，叫他在我为福音被锁链捆锁时，替我服事你。可是没有你的同意，我什么也不愿意做，好使你的善行不是出于勉强，而是出于甘心。同样，在\BibleBook{}{申命纪}中：\BibleQuote{我已将生命和死亡、善与恶，都摆在你面前……你要选择生命，好使你生存。}在\BibleBook{}{德训篇}中：\BibleQuote{天主从起初造了人，将他交在自己决断的手中；又为他颁布诫命和规定：你若愿意——在将来成就可接纳的信德，它们必拯救你。祂在你面前放置了火和水，你愿伸手向哪里，就伸向哪里。在人面前是善与恶、生命与死亡；贫穷与尊荣来自天主上主。}}\BibleRef{德 15:14}同样在\BibleBook{}{依撒意亚}中，我们读到：\BibleQuote{如果你们愿意，听从我，你们必享用地上的美物；但如果你们不愿意，不听从我，刀剑必吞灭你们：因为上主的口说了这话。}\BibleRef{依 1:19}如今他们虽百般掩饰，却在此暴露了他们的目的；因为他们公然企图驳斥天主的恩宠与仁慈，而当我们念经文时，正是想要获得这恩宠与仁慈：\BibleQuote{愿祢的旨意奉行在人间，如同在天上；}\BibleRef{玛 6:10}或这经文：\BibleQuote{不要让我们陷于诱惑，但救我们免于凶恶。}\BibleRef{玛 6:13}因为，如果结果是出于那愿意的和那奔跑的，而不是出于那施以仁慈的天主，我们为何要如此恳切地献上这些祈求呢？这并非说结果与我们的意愿无关，而是说我们的意愿不能成就结果，除非它得到天主神性的助佑。信德的健全性正在于此，它使我们\BibleQuote{寻求，就能找到；祈求，就能得到；敲门，门就给我们开。}\BibleRef{路 11:9}而反对这话的人，实在是将天主仁慈的门向自己关闭了。我不愿再多谈如此重要的事，因为我宁愿将它交托给信友的呻吟，而不是我自己的言辞。\par

% structure: p0103 source=h3 target=subsection reason=relative-heading-rank; base=h2

\subsection{(41) 伯拉纠式解经实例}

但我请求你们看看，他终究提出了什么样的反对意见：对于那\Quote{愿意并奔跑}的人，不需要天主的仁慈，而这仁慈实际上是先于他，好使他能够奔跑——因为，宗徒论到某人说：\BibleQuote{让他做他所愿意的}\BibleRef{格前 7:36}，——在我看来，就是他接着论述的事，当他说：\BibleQuote{他不犯罪，让他结婚吧！}\BibleRef{格前 7:36}仿佛愿意结婚应该被视为一件大事，而主题却是关于天主恩宠助佑的费力讨论，或者认为愿意结婚有任何重大益处，除非那掌管万有的天主眷顾将男女结合。或者，在宗徒\BibleBook{}{费肋孟书}的事例中，他说\Quote{你的善行不应像是出于勉强，而是出于甘心}——仿佛任何善行实际上是出于自愿，而非出于天主\BibleQuote{在我们内运行，使我有意愿也有能力，为成就祂的善意}\BibleRef{斐 2:13}。或者，当圣经在\BibleBook{}{申命纪}中说：\BibleQuote{祂将生命和死亡、善与恶，摆在人面前}，并劝诫他\BibleQuote{选择生命}；仿佛这劝诫本身不是来自天主的仁慈，或者选择生命有任何益处，除非天主激发爱德做出这样的选择，并在选择后赐予生命，关于这事有话说：\Quote{因为在他的震怒中有忿恼，在他的喜悦中有生命。}\par

或者，因为经上说：\BibleQuote{诫命，你若愿意，必拯救你}\BibleRef{德 15:15}——仿佛人不该感谢天主，因为他有意愿遵守诫命，因为如果他完全缺乏真理之光，就不可能拥有这样的意愿。\BibleQuote{火和水摆在人面前，人向自己所愿的方向伸手；}\BibleRef{德 15:16}然而，那召叫人走向更高召叫的，远超过人自己的任何思想，因为内心改正的开端在于信德，正如经上所写：\BibleQuote{你将从信德的开端而来，继续前进。}\BibleRef{歌 4:8}每个人都选择善，\BibleQuote{照天主所分给各人的信德尺度}\BibleRef{罗 12:3}；并且正如信德的元首所说：\BibleQuote{除非那派遣我的父吸引他，没有人能到我这里来。}\BibleRef{若 6:44}祂随后在说这话时充分清楚地解释了这是指相信祂的信德：\BibleQuote{我对你们所说的话，就是神，就是生命；但你们中间有不信的人。原来耶稣从起初就知道谁不信，谁要出卖祂。祂又说：为此我对你们说过：除非蒙我父的恩赐，没有人能到我这里来。}\BibleRef{若 6:62}\par

% structure: p0106 source=h3 target=subsection reason=relative-heading-rank; base=h2

\subsection{(42) 天主的应许是有条件的；旧约的圣人们是靠基督的恩宠得救的。}

然而，他以为在依撒意亚先知书中发现了支持其论点的有力依据，因为天主藉先知说：\Quote{如果你们乐意，且听从我，你们将享用地上的美物；但如果你们不乐意，且不听从我，刀剑必吞噬你们：因为上主亲口说了这话。} \BibleRef{依 1:19} 仿佛全部法律不都是充满这类条件；或者，它的诫命赐给骄傲的人，除了因为\Quote{法律是为了过犯而加添的，直到那蒙应许的后裔来到为止。} \BibleRef{迦 3:19} \Quote{因此，法律介入，为叫过犯增多；但罪恶在那里增多，恩宠在那里也格外丰富。} \BibleRef{罗 5:20} 换言之，使人领受诫命，却信赖自己的资源，并在这些资源中失败而成为犯法者，从而寻求解救者和救主；法律的恐惧能贬抑他，并如启蒙师一样引领他走向信德和恩宠。因此，\Quote{他们的软弱增多，他们就急忙追赶；} 为医治他们，基督在适当的时候来临。古代义人也是因祂的恩宠相信，并赖同一恩宠得助；他们欣然接受了关于祂的预知，甚至有些人预言了祂的来临——无论他们是在以色列民中，如梅瑟、农的儿子若苏厄、撒慕尔、达味及其他人；或是在那民族之外，如约伯；或是在那民族之前，如亚巴郎、诺厄及圣经中提及或未提及的一切其他人。\Quote{因为只有一位天主，在天主与人之间也只有一位中保，就是那人基督耶稣，} \BibleRef{弟前 2:5} 没有祂的恩宠，无人能脱离定罪，无论那定罪是从那众人都在其中犯了罪者承受的，还是后来因自己的罪愆加重的。\par

% structure: p0108 source=h2 target=section reason=relative-heading-rank; base=h2

\section{第二十章}

% structure: p0109 source=h3 target=subsection reason=relative-heading-rank; base=h2

\subsection{（43）无人得助，除非他自己也作工。我们的路程是持续的进步。}

但他最后陈述的含意是什么：\Quote{若有人说：\InnerQuote{人可能连在言语上也不犯罪吗？}那么，} 他说，\Quote{必须回答：\InnerQuote{完全可能，如果天主愿意；天主确实愿意，所以这是可能的。}} 看他多么不愿说：\Quote{如果天主赐助，那就可能；} 然而圣咏作者这样向天主呼求：\Quote{求祢作我的助佑，不要抛弃我；} 这里寻求的帮助当然不是为了获得肉身利益和避免肉身灾祸，而是为了实践和完成正义。因此我们说：\Quote{不要让我们陷入诱惑，但救我们脱离凶恶。} \BibleRef{玛 6:13} 没有人得助，除非他自己也作些什么；但他得助，如果他祈祷，如果他相信，如果他\Quote{照天主的旨意蒙召；} \BibleRef{罗 8:28} 因为\Quote{祂所预知的人，也预定他们与祂儿子的肖像相同，使祂在许多弟兄中作长子。不仅如此，祂所预定的人，也召叫了他们；祂所召叫的人，也使他们成义；祂所成义的人，也光荣了他们。} \BibleRef{罗 8:29} 因此，我们奔跑，每当我们前进；我们的健全与我们一起前进（正如伤口被说成是流淌的，当伤口在健康而谨慎的治疗过程中），以便我们在各方面完美，没有任何罪的软弱——这结果天主不仅愿意，而且甚至促成并帮助我们完成。而天主的恩宠与我们合作，通过我们的主耶稣基督，既藉着诫命、圣事和榜样，也藉着祂的圣神；藉着祂，那爱 \Quote{藉着圣神已倾注在我们心中} \BibleRef{罗 5:5} \Quote{以无可言喻的叹息为我们转求} \BibleRef{罗 8:26}，直到健全和救恩在我们里面完成，天主在祂永恒的真理中向我们显现，正如我们将看见祂。\par

% structure: p0111 source=h2 target=section reason=relative-heading-rank; base=h2

\section{第二十一章}

% structure: p0112 source=h3 target=subsection reason=relative-heading-rank; base=h2

\subsection{（44）本著作的结论。在重生者中，犯罪的是同意，而非私欲偏情。}

那么，凡认为任何人或任何一群人（除了天主与人之间的唯一中保\BibleRef{弟前 2:5}之外）曾在现今这个状态中活过或正在活着，而未曾需要或不需要罪的赦免，他就是反对圣经，因为宗徒在圣经中说：\BibleQuote{藉着一个人，罪恶进入了世界，藉着罪恶死亡，于是死亡传给了所有人，因为众人都犯了罪。}\BibleRef{罗 5:12} 他必然还要以不敬的争论断言，可能有人无需唯一中保基督的解放和救恩就能从罪恶中获得自由和得救。然而正是祂说过：\BibleQuote{健康的人不需要医生，有病的人才需要。}\BibleRef{玛 9:12} \BibleQuote{我来不是召义人，而是召罪人悔改。}\BibleRef{玛 9:13} 再者，凡说任何人在获得罪赦之后，在这个身体中曾生活或将生活得如此正义，以至于完全没有罪，他就是反驳宗徒若望，后者曾声明：\BibleQuote{如果我们说我们没有罪，就是欺骗自己，真理就不在我们内。}\BibleRef{若一 1:8} 请注意，说法不是\Emph{我们曾有}，而是\BibleQuote{\Emph{我们有}}。\newline{} 然而，如果有人争论说，宗徒的话涉及那按首位犯罪之人的意志所造成的缺陷而居于我们可朽肉身中的罪，宗徒保禄曾命令我们\BibleQuote{不可}顺从它的私欲而服从它，\BibleRef{罗 6:12} —— 那么，那完全不同意此内居之罪并因此不让它产生任何恶行（无论是行为、言语或思想）的人就不犯罪 —— 尽管对它的贪欲可能被激起（这贪欲在另一种意义上被称为罪，因为同意它就等于犯罪），但却是违背我们的意志被激起 —— 他一定是在作微妙的区分，并应考虑这一切与天主经有何关系，我们在其中说：\BibleQuote{宽免我们的罪债。}\BibleRef{玛 6:12} 现在，如果我判断正确，若我们从未在任何程度上同意上述罪的私欲，无论是言语失当还是放荡思想，就不必献上这样的祈祷；只需说：\BibleQuote{不要让我们陷于诱惑，但救我们免于凶恶。}\BibleRef{玛 6:13} 宗徒雅各伯也不可能说：\BibleQuote{我们众人在许多事上都有过犯。}\BibleRef{雅 3:2} 因为事实上，只有那被邪恶的贪欲所说服（无论是通过欺骗还是强迫）去做、说或想一些他本应避免之事的人，才犯罪，这贪欲引导他的欲望或厌憎违背正义的准则。最后，如果有人断言，在今世中曾经或现在有一些人（除了我们伟大的元首、\Quote{祂身体的救主}之外）毫无罪地是正义的 —— 这或者是由于不同意那些私欲，或者是因为这种罪不被视为罪，因天主因他们的虔诚生活而不归咎于他们（尽管无罪的福乐与罪不被归咎的福乐是两回事），—— 我认为不必过多争论此事。我很清楚有些人持此见解，我无勇气谴责他们的观点，但同时也不能为其辩护。但若有人说我们不应使用\BibleQuote{不要让我们陷于诱惑}的祈祷（凡主张为避免犯罪不需要天主的帮助，人的意志只须接受法律就足够的人，就是这样说的），那么我毫不犹豫地立即断言，这样的人应当被逐出公众的听闻，并被众口诅咒。\par

\SourceRecord{Translated by Peter Holmes and Robert Ernest Wallis, and revised by Benjamin B. Warfield.；Nicene and Post-Nicene Fathers, First Series；Vol. 5.；Edited by Philip Schaff.；Buffalo, NY: Christian Literature Publishing Co.,；1887.；<http://www.newadvent.org/fathers/1504.htm>.}
